\documentclass[11pt]{article}

\usepackage[margin=1in]{geometry}
\usepackage{amsmath,amssymb,amsthm,mathtools}
\usepackage{bm}
\usepackage{hyperref}

\newtheorem{theorem}{Theorem}
\newtheorem{proposition}{Proposition}
\newtheorem{example}{Example}
\newtheorem{corollary}{Corollary}
\newtheorem{lemma}{Lemma}
\newtheorem{remark}{Remark}

\DeclareMathOperator{\dist}{dist}
\newcommand{\R}{\mathbb{R}}
\newcommand{\ip}[2]{\left\langle #1,#2\right\rangle}
\newcommand{\norm}[1]{\left\lVert #1\right\rVert}

\title{Finite-Distortion Estimates and Modulus-Computed Level Sets}
\author{}
\date{}

\begin{document}

\maketitle

\section{Setup}

Let $\phi$ be a strictly convex barrier on a convex domain $\Omega\subseteq\R^n$. Define
\[
    g(x) := \nabla \phi(x),
    \qquad
    H(x) := \nabla^2 \phi(x).
\]
The Hessian metric induces the local primal and dual norms
\[
    \norm{u}_x^2 := u^\top H(x)u,
    \qquad
    \norm{w}_{x,*}^2 := w^\top H(x)^{-1}w.
\]
Let $d_R$ denote the Riemannian distance induced by the Hessian metric $H$.

For $x,y\in\Omega$, define the Jeffreys divergence
\[
    V(x,y)
    :=
    \ip{x-y}{g(x)-g(y)}.
\]
We also define the endpoint primal-dual residual
\[
    d(x,y)
    :=
    \norm{x-y}_x^2
    +
    \norm{g(x)-g(y)}_{x,*}^2.
\]

Throughout, assume global two-sided Hessian stability with modulus
\[
    \rho:[0,\infty)\to[1,\infty),
    \qquad
    \rho(0)=1,
    \qquad
    \rho \text{ nondecreasing},
\]
meaning that for all $x,z\in\Omega$,
\[
    \rho(d_R(x,z))^{-1}H(x)
    \preceq
    H(z)
    \preceq
    \rho(d_R(x,z))H(x).
\]

For $R\ge 0$, define
\[
    \Psi_\rho(R)
    :=
    2\int_0^R (R-r)\sqrt{\rho(r)}\,dr,
\]
and define the finite-distortion remainder
\[
    D_\rho(R)
    :=
    2\int_0^R
    (R-r)\left(\sqrt{\rho(r)}-1\right)\,dr.
\]
Thus
\[
    \Psi_\rho(R)=R^2+D_\rho(R).
\]

\section{The valid geodesic upper envelope}

\begin{proposition}[Upper envelope for the Jeffreys divergence]
For all $x,y\in\Omega$,
\[
    V(x,y)
    \le
    \Psi_\rho(d_R(x,y))
    =
    d_R(x,y)^2 + D_\rho(d_R(x,y)).
\]
\end{proposition}

\begin{proof}
Let $\gamma:[0,R]\to\Omega$ be a unit-speed geodesic from $y$ to $x$, where
$R=d_R(x,y)$. Then
\[
    x-y=\int_0^R \gamma'(a)\,da,
    \qquad
    g(x)-g(y)=\int_0^R H(\gamma(b))\gamma'(b)\,db.
\]
Thus
\[
    V(x,y)
    =
    \int_0^R\int_0^R
    \ip{\gamma'(a)}{H(\gamma(b))\gamma'(b)}\,da\,db.
\]
By Cauchy-Schwarz in the metric at $\gamma(b)$,
\[
    \ip{\gamma'(a)}{H(\gamma(b))\gamma'(b)}
    \le
    \norm{\gamma'(a)}_{\gamma(b)}\norm{\gamma'(b)}_{\gamma(b)}.
\]
Since $\gamma$ has unit speed, $\norm{\gamma'(b)}_{\gamma(b)}=1$. By Hessian stability,
\[
    H(\gamma(b))
    \preceq
    \rho(d_R(\gamma(a),\gamma(b)))H(\gamma(a))
    \preceq
    \rho(|a-b|)H(\gamma(a)).
\]
Hence
\[
    \norm{\gamma'(a)}_{\gamma(b)}
    \le
    \sqrt{\rho(|a-b|)}.
\]
Therefore
\[
    V(x,y)
    \le
    \int_0^R\int_0^R \sqrt{\rho(|a-b|)}\,da\,db.
\]
By symmetry of the square,
\[
    \int_0^R\int_0^R \sqrt{\rho(|a-b|)}\,da\,db
    =
    2\int_0^R (R-r)\sqrt{\rho(r)}\,dr.
\]
This proves the claim.
\end{proof}

\begin{proposition}[Endpoint finite-distortion identity]
Let $\gamma:[0,R]\to\Omega$ be a unit-speed geodesic and define
\[
    w(s):=V(\gamma(s),\gamma(0)).
\]
Then
\[
    w(R)=\int_0^R (R-s)w''(s)\,ds
\]
and
\[
    \int_0^R (R-s)w''(s)\,ds
    \le
    R^2+D_\rho(R).
\]
\end{proposition}

\begin{proof}
Since $w(0)=0$ and $w'(0)=0$, Taylor's formula gives
\[
    w(R)=\int_0^R (R-s)w''(s)\,ds.
\]
The previous proposition applied to the pair $(\gamma(R),\gamma(0))$ gives
\[
    w(R)=V(\gamma(R),\gamma(0))
    \le
    R^2+D_\rho(R).
\]
\end{proof}

\section{Target-dependent second variation}

The endpoint estimate above should not be confused with a target-independent estimate for
$V(x(t),y)$ when $y$ is an external fixed target. In that case a residual term appears.

Let $x(t)$ be a unit-time geodesic with constant kinetic energy
\[
    E:=\norm{\dot x(t)}_{x(t)}^2.
\]
For a unit-time geodesic, the step length is
\[
    R=d_R(x(0),x(1))=\sqrt{E}.
\]
Fix $y\in\Omega$ and define
\[
    v(t):=V(x(t),y).
\]
Define the target residual
\[
    U(t)
    :=
    (x(t)-y)
    -
    H(x(t))^{-1}\bigl(g(x(t))-g(y)\bigr).
\]

\begin{proposition}[Residual-corrected second variation]
Along the geodesic $x(t)$,
\[
    v''(t)
    =
    2E-\ip{\ddot x(t)}{U(t)}_{x(t)},
\]
where $\ip{a}{b}_{x}:=a^\top H(x)b$. Consequently,
\[
    v(1)
    =
    v(0)+v'(0)+E+\mathcal R_y[x],
\]
where
\[
    \mathcal R_y[x]
    :=
    -\int_0^1 (1-t)\ip{\ddot x(t)}{U(t)}_{x(t)}\,dt.
\]
In particular, if
\[
    v'(0)\le -(E+v(0)),
\]
then
\[
    V(x(1),y)=v(1)\le \mathcal R_y[x].
\]
\end{proposition}

\begin{proof}
Differentiate
\[
    v(t)=\ip{x(t)-y}{g(x(t))-g(y)}.
\]
The first derivative is
\[
    v'(t)
    =
    \ip{\dot x(t)}{g(x(t))-g(y)}
    +
    \ip{x(t)-y}{H(x(t))\dot x(t)}.
\]
Differentiating again gives
\[
\begin{aligned}
    v''(t)
    &=
    \ip{\ddot x(t)}{g(x(t))-g(y)}
    +
    \ip{\dot x(t)}{H(x(t))\dot x(t)} \\
    &\quad
    +
    \ip{\dot x(t)}{H(x(t))\dot x(t)}
    +
    D^3\phi(x(t))[\dot x(t),\dot x(t),x(t)-y]
    +
    \ip{x(t)-y}{H(x(t))\ddot x(t)}.
\end{aligned}
\]
The two middle quadratic terms equal $2E$. For a Hessian geodesic,
\[
    \ip{\ddot x(t)}{w}_{x(t)}
    =
    -\frac12 D^3\phi(x(t))[\dot x(t),\dot x(t),w]
\]
for every vector $w$. Therefore
\[
    D^3\phi(x(t))[\dot x(t),\dot x(t),x(t)-y]
    =
    -2\ip{\ddot x(t)}{x(t)-y}_{x(t)}.
\]
Substituting yields
\[
\begin{aligned}
    v''(t)
    &=
    2E
    +
    \ip{\ddot x(t)}{g(x(t))-g(y)}
    +
    \ip{x(t)-y}{H(x(t))\ddot x(t)}
    -
    2\ip{\ddot x(t)}{x(t)-y}_{x(t)} \\
    &=
    2E
    +
    \ip{\ddot x(t)}{g(x(t))-g(y)}
    -
    \ip{\ddot x(t)}{x(t)-y}_{x(t)} \\
    &=
    2E
    -
    \ip{\ddot x(t)}{(x(t)-y)-H(x(t))^{-1}(g(x(t))-g(y))}_{x(t)}.
\end{aligned}
\]
This is the claimed identity. Taylor's formula gives
\[
    v(1)=v(0)+v'(0)+\int_0^1(1-t)v''(t)\,dt.
\]
Using $\int_0^1(1-t)2E\,dt=E$ gives the residual-corrected formula.
\end{proof}


\section{A geodesic bound for the asymmetry residual}

The residual appearing in the target-dependent second variation has additional
structure.  For two points $x,y\in\Omega$, define
\[
    U(x,y)
    :=
    (x-y)-H(x)^{-1}(g(x)-g(y)).
\]
This is the same quantity as $U(t)$ above with $x=x(t)$.  It measures the
failure of the first-order Hessian approximation at $x$ to reproduce the
primal displacement from $x$ to $y$.  In the quadratic or constant-Hessian
case it vanishes identically.

Define the residual modulus
\[
    J_\rho(R)
    :=
    \int_0^R
    \left(
        \sqrt{\rho(r)}-\frac{1}{\sqrt{\rho(r)}}
    \right)dr .
\]

\begin{proposition}[Geodesic residual bound]
For all $x,y\in\Omega$,
\[
    \norm{U(x,y)}_x
    \le
    J_\rho(d_R(x,y)).
\]
Consequently, if $d_R(x,y)^2\le V(x,y)$, then
\[
    \norm{U(x,y)}_x
    \le
    J_\rho\!\left(\sqrt{V(x,y)}\right).
\]
\end{proposition}

\begin{proof}
Let $R=d_R(x,y)$, and let $\gamma:[0,R]\to\Omega$ be a unit-speed geodesic
with $\gamma(0)=x$ and $\gamma(R)=y$.  Then
\[
    x-y=-\int_0^R \gamma'(s)\,ds,
\]
and
\[
    g(x)-g(y)
    =
    -\int_0^R H(\gamma(s))\gamma'(s)\,ds.
\]
Thus
\[
\begin{aligned}
    U(x,y)
    &=(x-y)-H(x)^{-1}(g(x)-g(y)) \\
    &=\int_0^R
    \left[H(x)^{-1}H(\gamma(s))-I\right]\gamma'(s)\,ds .
\end{aligned}
\]
It remains to bound the norm of the integrand.  Set
\[
    A_s:=H(x)^{-1/2}H(\gamma(s))H(x)^{-1/2},
    \qquad
    z_s:=H(x)^{1/2}\gamma'(s).
\]
By Hessian stability and $d_R(x,\gamma(s))=s$,
\[
    \rho(s)^{-1}I\preceq A_s\preceq \rho(s)I.
\]
Moreover, since $\gamma$ has unit speed,
\[
    z_s^\top A_s z_s
    =
    \gamma'(s)^\top H(\gamma(s))\gamma'(s)
    =1.
\]
The $x$-norm of the integrand is
\[
\begin{aligned}
    \left\|
    \left[H(x)^{-1}H(\gamma(s))-I\right]\gamma'(s)
    \right\|_x
    &=
    \|(A_s-I)z_s\|_2 .
\end{aligned}
\]
Diagonalizing $A_s$, and using $z_s^\top A_s z_s=1$, gives
\[
    \|(A_s-I)z_s\|_2^2
    \le
    \max_{\lambda\in[\rho(s)^{-1},\rho(s)]}
    \frac{(\lambda-1)^2}{\lambda}.
\]
The maximum equals
\[
    \rho(s)+\rho(s)^{-1}-2
    =
    \left(\sqrt{\rho(s)}-\frac{1}{\sqrt{\rho(s)}}\right)^2.
\]
Therefore
\[
    \left\|
    \left[H(x)^{-1}H(\gamma(s))-I\right]\gamma'(s)
    \right\|_x
    \le
    \sqrt{\rho(s)}-\frac{1}{\sqrt{\rho(s)}}.
\]
Integrating over $s\in[0,R]$ yields
\[
    \norm{U(x,y)}_x
    \le
    \int_0^R
    \left(
        \sqrt{\rho(s)}-\frac{1}{\sqrt{\rho(s)}}
    \right)ds
    =J_\rho(R).
\]
The final claim follows from monotonicity of $J_\rho$ and the coercivity
$d_R(x,y)^2\le V(x,y)$.
\end{proof}

\begin{corollary}[Exponential stability specialization]
Suppose
\[
    \rho(r)=e^{Lr}.
\]
Then
\[
    J_\rho(R)
    =
    \int_0^R \left(e^{Lr/2}-e^{-Lr/2}\right)dr
    =
    \frac{4}{L}\left(\cosh\left(\frac{LR}{2}\right)-1\right).
\]
Therefore
\[
    \norm{U(x,y)}_x
    \le
    \frac{4}{L}
    \left(
        \cosh\left(\frac{L}{2}d_R(x,y)\right)-1
    \right).
\]
For the standard self-concordant normalization $L=2$, this becomes
\[
    \norm{U(x,y)}_x
    \le
    2\left(\cosh(d_R(x,y))-1\right).
\]
In particular, as $R=d_R(x,y)\downarrow0$,
\[
    J_\rho(R)=\frac{L}{2}R^2+O(R^4),
\]
so the residual is second order in geodesic distance.
\end{corollary}



\section{A corrected convergence statement}

\begin{theorem}[Residual-controlled one-step convergence]
Assume global two-sided Hessian stability with modulus $\rho$. Let $T(x)=x(1)$ be obtained from a unit-time geodesic $x(t)$ starting at $x(0)=x$, with step length $R=d_R(x,T(x))$ and kinetic energy $E=R^2$. Fix $y\in\Omega$ and define $v(t)=V(x(t),y)$.

If
\[
    v'(0)\le -(E+v(0)),
\]
then
\[
    V(T(x),y)\le \mathcal R_y[x],
\]
where
\[
    \mathcal R_y[x]
    :=
    -\int_0^1 (1-t)\ip{\ddot x(t)}{U(t)}_{x(t)}\,dt,
\]
and
\[
    U(t)
    :=
    (x(t)-y)-H(x(t))^{-1}\bigl(g(x(t))-g(y)\bigr).
\]
Consequently, if a scalar function $G$ satisfies
\[
    \mathcal R_y[x]\le G(V(x,y))
\]
for all $x$ in a level set $\mathcal L_\delta(y):=\{x:V(x,y)\le\delta\}$, then
\[
    V(T(x),y)\le G(V(x,y))
    \qquad
    \forall x\in\mathcal L_\delta(y).
\]
If $G(v)\le v$ for all $v\in[0,\delta]$, then $\mathcal L_\delta(y)$ is invariant. If $G(v)\le v^2$ for all $v\in[0,\delta]$, then the iteration $x_{k+1}=T(x_k)$ satisfies
\[
    V(x_{k+1},y)\le V(x_k,y)^2
\]
whenever $x_k\in\mathcal L_\delta(y)$.
\end{theorem}

\begin{proof}
This is immediate from the residual-corrected second variation formula. Under the descent condition,
\[
    V(T(x),y)=v(1)
    \le
    \mathcal R_y[x].
\]
The conclusions follow by substituting the scalar upper bound $G$ and iterating the scalar inequality.
\end{proof}




\section{A modulus-computed one-step map}

The previous theorem is correct but not yet useful unless the residual term is bounded.
We now show that, under the same update hypotheses and using only the stability modulus,
one obtains a computable scalar map for the target divergence.

Define
\[
    L_\rho
    :=
    \limsup_{r\downarrow0}\frac{\log \rho(r)}{r},
\]
and assume $L_\rho<\infty$.  Also define the first-order stability primitive
\[
    S_\rho(R)
    :=
    \int_0^R \sqrt{\rho(r)}\,dr.
\]

\begin{lemma}[Endpoint norm bounds from the modulus]
For all $a,b\in\Omega$, with $R=d_R(a,b)$,
\[
    \norm{a-b}_a
    \le
    S_\rho(R),
    \qquad
    \norm{g(a)-g(b)}_{a,*}
    \le
    S_\rho(R).
\]
Consequently,
\[
    d(a,b)
    \le
    2S_\rho(d_R(a,b))^2.
\]
\end{lemma}

\begin{proof}
Let $\gamma:[0,R]\to\Omega$ be a unit-speed geodesic with $\gamma(0)=a$ and
$\gamma(R)=b$.  Then
\[
    b-a=\int_0^R \gamma'(s)\,ds.
\]
Using Hessian stability between $a$ and $\gamma(s)$,
\[
    H(a)\preceq \rho(s)H(\gamma(s)).
\]
Therefore
\[
    \norm{\gamma'(s)}_a
    \le
    \sqrt{\rho(s)}\norm{\gamma'(s)}_{\gamma(s)}
    =
    \sqrt{\rho(s)}.
\]
Thus
\[
    \norm{b-a}_a
    \le
    \int_0^R \sqrt{\rho(s)}\,ds
    =
    S_\rho(R).
\]
For the dual term,
\[
    g(b)-g(a)=\int_0^R H(\gamma(s))\gamma'(s)\,ds.
\]
By duality,
\[
\begin{aligned}
    \norm{g(b)-g(a)}_{a,*}
    &=
    \sup_{\norm{w}_a\le1}
    \int_0^R \ip{w}{H(\gamma(s))\gamma'(s)}\,ds \\
    &\le
    \sup_{\norm{w}_a\le1}
    \int_0^R \norm{w}_{\gamma(s)}\norm{\gamma'(s)}_{\gamma(s)}\,ds.
\end{aligned}
\]
Again by stability, $H(\gamma(s))\preceq \rho(s)H(a)$, so
$\norm{w}_{\gamma(s)}\le\sqrt{\rho(s)}\norm{w}_a$.  Since the geodesic has unit speed,
\[
    \norm{g(b)-g(a)}_{a,*}
    \le
    \int_0^R \sqrt{\rho(s)}\,ds
    =
    S_\rho(R).
\]
Squaring and adding gives the bound on $d(a,b)$.
\end{proof}

\begin{lemma}[Infinitesimal acceleration bound]
Let $x(t)$ be a Hessian geodesic with constant kinetic energy
\[
    E=\norm{\dot x(t)}_{x(t)}^2.
\]
Then
\[
    \norm{\ddot x(t)}_{x(t)}
    \le
    \frac{L_\rho}{2}E.
\]
\end{lemma}

\begin{proof}
The local limit of the stability condition gives, for all $x,u,v$,
\[
    |D^3\phi(x)[v,u,u]|
    \le
    L_\rho\norm{v}_x\norm{u}_x^2.
\]
Indeed, apply the upper and lower Hessian stability inequalities to
$x+\epsilon v/\norm{v}_x$ and let $\epsilon\downarrow0$.
For a Hessian geodesic,
\[
    \ip{\ddot x(t)}{w}_{x(t)}
    =
    -\frac12 D^3\phi(x(t))[\dot x(t),\dot x(t),w].
\]
Hence, for $\norm{w}_{x(t)}\le1$,
\[
    |\ip{\ddot x(t)}{w}_{x(t)}|
    \le
    \frac{L_\rho}{2}\norm{\dot x(t)}_{x(t)}^2
    =
    \frac{L_\rho}{2}E.
\]
Taking the supremum over $\norm{w}_{x(t)}\le1$ gives the claim.
\end{proof}

For $v\ge0$, define
\[
    \Lambda_\rho(v)
    :=
    \sqrt v + \sqrt{2}\,S_\rho(\sqrt v),
\]
and define the modulus-computed one-step map
\[
    F_\rho(v)
    :=
    L_\rho\,S_\rho(\sqrt v)^2\,S_\rho(\Lambda_\rho(v)).
\]

\begin{theorem}[Modulus-computed one-step bound]
Assume global two-sided Hessian stability with modulus $\rho$, $L_\rho<\infty$, and
$ d_R(a,b)^2\le V(a,b)$ for all $a,b$.  Let $T(x)=x(1)$ be obtained from a
unit-time geodesic $x(t)$ starting at $x(0)=x$, with step length
$R=d_R(x,T(x))$ and kinetic energy $E=R^2$.  Fix $y\in\Omega$ and define
$v(t)=V(x(t),y)$.

Assume the update satisfies
\[
    v'(0)\le -(E+V(x,y))
\]
and the step compatibility condition
\[
    R^2\le d(x,y).
\]
Then
\[
    V(T(x),y)
    \le
    F_\rho(V(x,y)).
\]
Thus the scalar one-step estimate is derived from the modulus and the update
conditions; it is not an additional assumption.
\end{theorem}

\begin{proof}
Let $v_0:=V(x,y)$ and $r_0:=d_R(x,y)$.  By geodesic coercivity,
\[
    r_0\le \sqrt{v_0}.
\]
By the endpoint norm lemma,
\[
    d(x,y)
    \le
    2S_\rho(r_0)^2
    \le
    2S_\rho(\sqrt{v_0})^2.
\]
The step compatibility condition gives
\[
    R^2=E\le d(x,y)
    \le
    2S_\rho(\sqrt{v_0})^2.
\]
Therefore
\[
    R\le \sqrt{2}\,S_\rho(\sqrt{v_0}).
\]
For every $t\in[0,1]$,
\[
    d_R(x(t),y)
    \le
    d_R(x(t),x)+d_R(x,y)
    \le
    R+r_0
    \le
    \sqrt{2}\,S_\rho(\sqrt{v_0})+\sqrt{v_0}
    =
    \Lambda_\rho(v_0).
\]
Applying the endpoint norm lemma at the point $x(t)$ gives
\[
    \norm{x(t)-y}_{x(t)}
    \le
    S_\rho(\Lambda_\rho(v_0)),
\]
and
\[
    \norm{g(x(t))-g(y)}_{x(t),*}
    \le
    S_\rho(\Lambda_\rho(v_0)).
\]
Hence the target residual
\[
    U(t)=(x(t)-y)-H(x(t))^{-1}(g(x(t))-g(y))
\]
satisfies
\[
    \norm{U(t)}_{x(t)}
    \le
    2S_\rho(\Lambda_\rho(v_0)).
\]
By the residual-corrected second variation and the acceleration bound,
\[
\begin{aligned}
    \mathcal R_y[x]
    &:=
    -\int_0^1(1-t)\ip{\ddot x(t)}{U(t)}_{x(t)}\,dt \\
    &\le
    \int_0^1(1-t)
    \norm{\ddot x(t)}_{x(t)}\norm{U(t)}_{x(t)}\,dt \\
    &\le
    \int_0^1(1-t)
    \frac{L_\rho}{2}E\cdot 2S_\rho(\Lambda_\rho(v_0))\,dt \\
    &=
    \frac{L_\rho}{2}E S_\rho(\Lambda_\rho(v_0)).
\end{aligned}
\]
Using $E\le 2S_\rho(\sqrt{v_0})^2$ gives
\[
    \mathcal R_y[x]
    \le
    L_\rho S_\rho(\sqrt{v_0})^2S_\rho(\Lambda_\rho(v_0))
    =
    F_\rho(v_0).
\]
Finally, the descent condition and the residual-corrected formula give
\[
    V(T(x),y)\le \mathcal R_y[x].
\]
Combining the last two inequalities proves
\[
    V(T(x),y)\le F_\rho(V(x,y)).
\]
\end{proof}

\section{Level sets and rates from the computed scalar map}

\begin{corollary}[Modulus-computed invariant level set]
Let $F_\rho$ be the scalar map defined above.  Define
\[
    \delta_{\rho,\mathrm{inv}}
    :=
    \sup\left\{
        \delta\in(0,1]:
        F_\rho(v)\le v
        \ \text{for all }v\in[0,\delta]
    \right\}.
\]
If every update satisfies the hypotheses of the modulus-computed one-step theorem and
$V(x_0,y)\le\delta_{\rho,\mathrm{inv}}$, then the iterates $x_{k+1}=T(x_k)$ satisfy
\[
    V(x_k,y)\le\delta_{\rho,\mathrm{inv}}
    \qquad\text{for all }k.
\]
Moreover,
\[
    V(x_{k+1},y)
    \le
    F_\rho(V(x_k,y))
    \qquad\text{for all }k.
\]
\end{corollary}

\begin{proof}
If $V(x_k,y)\le\delta_{\rho,\mathrm{inv}}$, then the one-step theorem gives
\[
    V(x_{k+1},y)
    \le
    F_\rho(V(x_k,y))
    \le
    V(x_k,y)
    \le
    \delta_{\rho,\mathrm{inv}}.
\]
Induction proves the result.
\end{proof}

\begin{corollary}[Computed linear and superlinear rates]
For any $\delta\le\delta_{\rho,\mathrm{inv}}$, define
\[
    q_\rho(\delta)
    :=
    \sup_{0<v\le\delta}\frac{F_\rho(v)}{v}.
\]
Then every sequence initialized in $\mathcal L_\delta(y):=\{x:V(x,y)\le\delta\}$ satisfies
\[
    V(x_{k+1},y)
    \le
    q_\rho(\delta)V(x_k,y).
\]
In particular, if $q_\rho(\delta)<1$, the convergence is linearly contracting on
$\mathcal L_\delta(y)$ with contraction factor $q_\rho(\delta)$.

If $\sqrt{\rho(r)}=1+O(r)$ as $r\downarrow0$, then
\[
    F_\rho(v)=O(v^{3/2}),
\]
so $q_\rho(\delta)\to0$ as $\delta\downarrow0$.  Thus the computed level sets yield
asymptotically superlinear contraction.
\end{corollary}

\begin{proof}
The first claim follows by dividing the one-step bound
$V(x_{k+1},y)\le F_\rho(V(x_k,y))$ by $V(x_k,y)$ and taking the supremum over
$0<v\le\delta$.
For the asymptotic statement, $S_\rho(R)=R+O(R^2)$, hence
\[
    \Lambda_\rho(v)=(1+\sqrt2)\sqrt v+O(v).
\]
Therefore
\[
    F_\rho(v)
    =
    L_\rho\bigl(v+O(v^{3/2})\bigr)\bigl((1+\sqrt2)\sqrt v+O(v)\bigr)
    =
    O(v^{3/2}).
\]
\end{proof}

\begin{remark}[Quadratic contraction]
A quadratic level can still be defined by
\[
    \delta_{\rho,\mathrm{quad}}
    :=
    \sup\left\{
        \delta\in(0,1]:
        F_\rho(v)\le v^2
        \ \text{for all }v\in[0,\delta]
    \right\}.
\]
If this radius is positive, then initialization in $\mathcal L_{\delta_{\rho,\mathrm{quad}}}(y)$
gives $V(x_{k+1},y)\le V(x_k,y)^2$ by the same induction.  For the usual exponential
self-concordant modulus, the computable residual bound above scales like $v^{3/2}$;
therefore this particular modulus-only argument gives superlinear contraction.  A
quadratic theorem requires either a sharper algorithmic step estimate or a stronger
step scaling than $R^2\le d(x,y)$.
\end{remark}

\begin{corollary}[Distance initialization]
Let
\[
    r_{\rho,\mathrm{inv}}
    :=
    \sup\{r\ge0:\Psi_\rho(r)\le \delta_{\rho,\mathrm{inv}}\}.
\]
If $d_R(x_0,y)\le r_{\rho,\mathrm{inv}}$, then $V(x_0,y)\le\delta_{\rho,\mathrm{inv}}$ and
therefore the invariant-level-set conclusion holds.
\end{corollary}

\begin{proof}
By the geodesic upper envelope,
\[
    V(x_0,y)\le\Psi_\rho(d_R(x_0,y))\le\delta_{\rho,\mathrm{inv}}.
\]
Then apply the invariant level-set corollary.
\end{proof}

\section{Examples}

\begin{example}[Constant Hessian]
Suppose $H(x)\equiv H_0$ is constant. Then $\rho(r)\equiv1$, so
\[
    D_\rho(R)=0,
    \qquad
    L_\rho=0,
    \qquad
    S_\rho(R)=R.
\]
Consequently the modulus-computed one-step map is
\[
    F_\rho(v)=0.
\]
Geodesics are affine, so $\ddot x(t)=0$ and therefore
\[
    \mathcal R_y[x]=0
\]
for every target $y$. If the descent condition
\[
    v'(0)\le -(R^2+v(0))
\]
holds, then
\[
    V(T(x),y)=0.
\]
Thus the flat case has exact cancellation: the target-dependent residual term vanishes because the geodesic acceleration is zero.
\end{example}

\begin{example}[Exponential stability / self-concordant case]
Suppose the Hessian metric satisfies exponential stability
\[
    e^{-L d_R(x,z)}H(x)
    \preceq
    H(z)
    \preceq
    e^{L d_R(x,z)}H(x).
\]
Then
\[
    \rho(r)=e^{Lr},
    \qquad
    \sqrt{\rho(r)}=e^{Lr/2}.
\]
The geodesic upper envelope is
\[
    \Psi_\rho(R)
    =
    2\int_0^R (R-r)e^{Lr/2}\,dr
    =
    \frac{8}{L^2}\left(e^{LR/2}-1\right)-\frac{4}{L}R.
\]
The finite endpoint distortion is
\[
    D_\rho(R)
    =
    \frac{8}{L^2}\left(e^{LR/2}-1\right)-\frac{4}{L}R-R^2.
\]
For small $R$,
\[
    D_\rho(R)=\frac{L}{3}R^3+O(R^4).
\]
The modulus-computed residual map uses
\[
    L_\rho=L,
    \qquad
    S_\rho(R)=\frac{2}{L}\left(e^{LR/2}-1\right).
\]
Therefore
\[
    \Lambda_\rho(v)
    =
    \sqrt v+\frac{2\sqrt2}{L}\left(e^{L\sqrt v/2}-1\right),
\]
and
\[
    F_\rho(v)
    =
    L\left[\frac{2}{L}\left(e^{L\sqrt v/2}-1\right)\right]^2
    \left[\frac{2}{L}\left(e^{L\Lambda_\rho(v)/2}-1\right)\right].
\]
As $v\downarrow0$,
\[
    F_\rho(v)=L(1+\sqrt2)v^{3/2}+O(v^2).
\]
This gives a fully computable superlinear level-set map.  The endpoint distortion
$D_\rho$ controls $V(x(1),x(0))$; the target-dependent map $F_\rho$ controls
$V(x(1),y)$ under the descent and step compatibility hypotheses of the theorem.
\end{example}

\begin{example}[Self-concordant normalization $L=1$]
For $L=1$, the exponential modulus is
\[
    \rho(r)=e^r.
\]
The upper envelope becomes
\[
    \Psi_\rho(R)
    =
    8\left(e^{R/2}-1\right)-4R,
\]
and the finite endpoint distortion is
\[
    D_\rho(R)
    =
    8\left(e^{R/2}-1\right)-4R-R^2.
\]
As $R\downarrow0$,
\[
    D_\rho(R)=\frac13 R^3+O(R^4).
\]
Here
\[
    S_\rho(R)=2(e^{R/2}-1),
\]
so
\[
    \Lambda_\rho(v)
    =
    \sqrt v+2\sqrt2(e^{\sqrt v/2}-1),
\]
and
\[
    F_\rho(v)
    =
    \bigl[2(e^{\sqrt v/2}-1)\bigr]^2
    \bigl[2(e^{\Lambda_\rho(v)/2}-1)\bigr].
\]
As $v\downarrow0$,
\[
    F_\rho(v)=(1+\sqrt2)v^{3/2}+O(v^2).
\]
\end{example}

\begin{example}[Why no $\rho$-only target bound can hold]
Let $\phi(x)=-\log x$ on $(0,\infty)$. Then
\[
    H(x)=x^{-2},
    \qquad
    d_R(x,z)=|\log x-\log z|,
\]
and the metric has exponential stability. Let
\[
    x(t)=e^{Rt},
    \qquad
    x(0)=1,
\]
which is a geodesic of length $R$. For a fixed target $y>0$,
\[
    V(x(t),y)
    =
    \frac{e^{Rt}}{y}+y e^{-Rt}-2.
\]
Hence
\[
    V(x(1),y)-V(x(0),y)-v'(0)
    =
    \frac{e^R-1-R}{y}+y(e^{-R}-1+R).
\]
For fixed $R>0$, this expression is unbounded as $y\to\infty$. Therefore it cannot be bounded by a function of $R$ and $\rho$ alone. This demonstrates why the target residual $\mathcal R_y[x]$ is necessary.
\end{example}






\section{A self-concordant level-set estimate}

Let \(\phi\) be a standard self-concordant barrier with
\[
    g(x):=\nabla \phi(x),
    \qquad
    H(x):=\nabla^2\phi(x).
\]
For a vector \(u\) and covector \(w\), write
\[
    \|u\|_x^2:=u^\top H(x)u,
    \qquad
    \|w\|_{x,*}^2:=w^\top H(x)^{-1}w.
\]
Define the Jeffreys divergence
\[
    V(x,y):=\langle x-y,g(x)-g(y)\rangle .
\]

Let \(x(t)\), \(t\in[0,1]\), be a unit-time geodesic for the Hessian metric, with constant kinetic energy
\[
    E:=\|\dot x(t)\|_{x(t)}^2.
\]
Fix a target point \(y\), and define
\[
    v(t):=V(x(t),y),
    \qquad
    v_0:=v(0)=V(x(0),y).
\]
Define the asymmetry residual
\[
    U(x,y):=(x-y)-H(x)^{-1}(g(x)-g(y)).
\]

Assume the descent condition
\[
    v'(0)\le -(E+v_0).
\]
Assume also the energy compatibility condition
\[
    E\le d(x(0),y),
\]
where
\[
    d(x,y)
    :=
    \|x-y\|_x^2
    +
    \|g(x)-g(y)\|_{x,*}^2.
\]
Finally, assume the geodesic path remains in the current \(V\)-sublevel set:
\[
    v(t)\le v_0
    \qquad
    \forall t\in[0,1].
\]

Then
\[
    v(1)
    \le
    \frac{(e^{\sqrt{v_0}}-1)^4}{e^{\sqrt{v_0}}}.
\]
In particular,
\[
    v(1)=v_0^2+O(v_0^3)
    \qquad
    (v_0\downarrow0).
\]

\begin{proof}
For a self-concordant Hessian metric, the second-variation identity is
\[
    v''(t)
    =
    2E-\langle \ddot x(t),U(x(t),y)\rangle_{x(t)}.
\]
Taylor's formula gives the exact identity
\[
    v(1)
    =
    v_0+v'(0)+\int_0^1(1-t)v''(t)\,dt.
\]
Substituting the expression for \(v''\),
\[
    \int_0^1(1-t)v''(t)\,dt
    =
    E-\int_0^1(1-t)
    \langle \ddot x(t),U(x(t),y)\rangle_{x(t)}\,dt.
\]
Hence
\[
    v(1)
    =
    v_0+v'(0)+E
    -
    \int_0^1(1-t)
    \langle \ddot x(t),U(x(t),y)\rangle_{x(t)}\,dt.
\]
The descent condition gives
\[
    v_0+v'(0)+E\le 0.
\]
Therefore
\[
    v(1)
    \le
    \int_0^1(1-t)
    \left|
    \langle \ddot x(t),U(x(t),y)\rangle_{x(t)}
    \right|dt.
\]

By self-concordance and the geodesic equation,
\[
    \|\ddot x(t)\|_{x(t)}\le E.
\]
Thus
\[
    \left|
    \langle \ddot x(t),U(x(t),y)\rangle_{x(t)}
    \right|
    \le
    E\|U(x(t),y)\|_{x(t)}.
\]

The exponential geodesic residual bound for standard self-concordance is
\[
    \|U(z,y)\|_z
    \le
    2\left(\cosh\sqrt{V(z,y)}-1\right).
\]
Since \(v(t)=V(x(t),y)\le v_0\), this gives
\[
    \|U(x(t),y)\|_{x(t)}
    \le
    2\left(\cosh\sqrt{v_0}-1\right).
\]
Therefore
\[
\begin{aligned}
    v(1)
    &\le
    \int_0^1(1-t)
    E\cdot
    2\left(\cosh\sqrt{v_0}-1\right)dt  \\
    &=
    E\left(\cosh\sqrt{v_0}-1\right).
\end{aligned}
\]

It remains to bound \(E\) in terms of \(v_0\). By assumption,
\[
    E\le d(x(0),y).
\]
For standard self-concordance, exponential Hessian stability gives
\[
    d(x,y)
    \le
    2\left(e^{\sqrt{V(x,y)}}-1\right)^2.
\]
Thus
\[
    E
    \le
    2\left(e^{\sqrt{v_0}}-1\right)^2.
\]
Substituting this into the previous estimate yields
\[
    v(1)
    \le
    2\left(e^{\sqrt{v_0}}-1\right)^2
    \left(\cosh\sqrt{v_0}-1\right).
\]
Using
\[
    2(\cosh s-1)
    =
    \frac{(e^s-1)^2}{e^s},
\]
with \(s=\sqrt{v_0}\), we obtain
\[
    v(1)
    \le
    \frac{(e^{\sqrt{v_0}}-1)^4}{e^{\sqrt{v_0}}}.
\]
\end{proof}






\subsection*{Quadratic convergence interpretation}

Define
\[
    F(v)
    :=
    \frac{(e^{\sqrt v}-1)^4}{e^{\sqrt v}}.
\]
Then the estimate above says
\[
    v(1)\le F(v_0).
\]
Moreover,
\[
    \frac{F(v)}{v^2}
    =
    \left(
        \frac{e^{\sqrt v}-1}{\sqrt v}
    \right)^4
    e^{-\sqrt v}.
\]
Hence
\[
    \lim_{v\downarrow0}\frac{F(v)}{v^2}=1.
\]
Thus
\[
    v(1)\le v_0^2(1+o(1)).
\]
Equivalently, for every \(\varepsilon>0\), there exists \(\delta_\varepsilon>0\) such that
\[
    0<v_0\le \delta_\varepsilon
    \quad\Longrightarrow\quad
    v(1)\le (1+\varepsilon)v_0^2.
\]

For a finite-region statement, define
\[
    C_\delta
    :=
    \sup_{0<v\le\delta}\frac{F(v)}{v^2}
    =
    \left(
        \frac{e^{\sqrt\delta}-1}{\sqrt\delta}
    \right)^4
    e^{-\sqrt\delta}.
\]
Then, whenever \(v_0\le\delta\),
\[
    v(1)\le C_\delta v_0^2.
\]
For example,
\[
    C_{1/4}
    =
    \left(
        \frac{e^{1/2}-1}{1/2}
    \right)^4e^{-1/2}
    \approx 1.730.
\]
Therefore
\[
    v_0\le \frac14
    \quad\Longrightarrow\quad
    v(1)\le 1.730\,v_0^2.
\]

The same scalar bound gives strict contraction on the region where
\[
    F(v)<v.
\]
Numerically,
\[
    F(v)<v
    \qquad\text{for}\qquad
    0<v<0.4671709662\ldots.
\]
Thus this estimate certifies strict decrease for approximately
\[
    v_0<0.46717.
\]

\section{General modulus version}

Now suppose instead that the Hessian metric satisfies global two-sided stability with modulus
\[
    \rho:[0,\infty)\to[1,\infty),
    \qquad
    \rho(0)=1,
    \qquad
    \rho \text{ nondecreasing},
\]
meaning
\[
    \rho(d_R(x,z))^{-1}H(x)
    \preceq
    H(z)
    \preceq
    \rho(d_R(x,z))H(x)
    \qquad
    \forall x,z.
\]
Define
\[
    K_\rho(r):=\sqrt{\rho(r)},
\]
\[
    S_\rho(R):=\int_0^R K_\rho(s)\,ds,
\]
and
\[
    J_\rho(R):=
    \int_0^R
    \left(
        K_\rho(s)-K_\rho(s)^{-1}
    \right)ds.
\]
Assume also that the modulus has finite infinitesimal logarithmic slope
\[
    L_\rho
    :=
    \limsup_{r\downarrow0}
    \frac{\log\rho(r)}{r}
    <\infty.
\]

Let \(x(t)\), \(v(t)\), \(v_0\), \(E\), and \(U\) be as above. Assume
\[
    v'(0)\le -(E+v_0),
\]
\[
    E\le d(x(0),y),
\]
and
\[
    v(t)\le v_0
    \qquad
    \forall t\in[0,1].
\]
Then
\[
    v(1)
    \le
    F_\rho(v_0),
\]
where
\[
    F_\rho(v)
    :=
    \frac{L_\rho}{2}
    S_\rho(\sqrt v)^2
    J_\rho(\sqrt v).
\]

\begin{proof}
The same second-variation identity gives
\[
    v(1)
    \le
    \int_0^1(1-t)
    \left|
    \langle \ddot x(t),U(x(t),y)\rangle_{x(t)}
    \right|dt.
\]

The finite infinitesimal slope \(L_\rho\) gives the geodesic acceleration bound
\[
    \|\ddot x(t)\|_{x(t)}
    \le
    \frac{L_\rho}{2}E.
\]
Hence
\[
    \left|
    \langle \ddot x(t),U(x(t),y)\rangle_{x(t)}
    \right|
    \le
    \frac{L_\rho}{2}E\|U(x(t),y)\|_{x(t)}.
\]

The modulus gives the residual bound
\[
    \|U(z,y)\|_z
    \le
    J_\rho(d_R(z,y)).
\]
Since \(d_R(z,y)^2\le V(z,y)\), and since \(v(t)\le v_0\), we get
\[
    \|U(x(t),y)\|_{x(t)}
    \le
    J_\rho(\sqrt{v_0}).
\]
Therefore
\[
    v(1)
    \le
    \frac{L_\rho}{2}EJ_\rho(\sqrt{v_0})
    \int_0^1(1-t)dt
    =
    \frac{L_\rho}{4}EJ_\rho(\sqrt{v_0}).
\]

It remains to bound \(E\). By the energy compatibility condition,
\[
    E\le d(x(0),y).
\]
Let \(R=d_R(x(0),y)\). The modulus gives
\[
    \|x(0)-y\|_{x(0)}\le S_\rho(R),
\]
and
\[
    \|g(x(0))-g(y)\|_{x(0),*}\le S_\rho(R).
\]
Therefore
\[
    d(x(0),y)
    \le
    2S_\rho(R)^2.
\]
Since
\[
    R^2=d_R(x(0),y)^2\le V(x(0),y)=v_0,
\]
we have
\[
    E
    \le
    2S_\rho(\sqrt{v_0})^2.
\]
Substituting this into the previous bound gives
\[
    v(1)
    \le
    \frac{L_\rho}{4}
    \left[
        2S_\rho(\sqrt{v_0})^2
    \right]
    J_\rho(\sqrt{v_0})
    =
    \frac{L_\rho}{2}
    S_\rho(\sqrt{v_0})^2
    J_\rho(\sqrt{v_0}).
\]
\end{proof}

\subsection*{Quadratic rate for a general modulus}

If
\[
    \rho(r)=1+L_\rho r+O(r^2)
    \qquad
    (r\downarrow0),
\]
then
\[
    S_\rho(\sqrt v)
    =
    \sqrt v+O(v),
\]
and
\[
    J_\rho(\sqrt v)
    =
    \frac{L_\rho}{2}v+O(v^{3/2}).
\]
Therefore
\[
    F_\rho(v)
    =
    \frac{L_\rho}{2}
    \left(v+O(v^{3/2})\right)
    \left(\frac{L_\rho}{2}v+O(v^{3/2})\right),
\]
so
\[
    F_\rho(v)
    =
    \frac{L_\rho^2}{4}v^2+O(v^{5/2}).
\]
Thus the generalized estimate is quadratically convergent in the sense that
\[
    v(1)\le C_\delta(\rho)v_0^2
\]
on any level set \(0<v_0\le\delta\) with
\[
    C_\delta(\rho)
    :=
    \sup_{0<v\le\delta}
    \frac{F_\rho(v)}{v^2}
    <\infty.
\]
Moreover,
\[
    \lim_{\delta\downarrow0}C_\delta(\rho)
    =
    \frac{L_\rho^2}{4}.
\]

\subsection*{Exponential modulus}

If
\[
    \rho(r)=e^{Lr},
\]
then
\[
    K_\rho(r)=e^{Lr/2},
\]
\[
    S_\rho(R)
    =
    \frac{2}{L}
    \left(e^{LR/2}-1\right),
\]
and
\[
    J_\rho(R)
    =
    \frac{4}{L}
    \left(
        \cosh\left(\frac{LR}{2}\right)-1
    \right).
\]
Hence
\[
\begin{aligned}
    F_L(v)
    &=
    \frac{L}{2}
    \left[
        \frac{2}{L}
        \left(e^{L\sqrt v/2}-1\right)
    \right]^2
    \cdot
    \frac{4}{L}
    \left[
        \cosh\left(\frac{L\sqrt v}{2}\right)-1
    \right]  \\
    &=
    \frac{8}{L^2}
    \left(e^{L\sqrt v/2}-1\right)^2
    \left[
        \cosh\left(\frac{L\sqrt v}{2}\right)-1
    \right].
\end{aligned}
\]
Using
\[
    2(\cosh a-1)=\frac{(e^a-1)^2}{e^a},
\]
this simplifies to
\[
    F_L(v)
    =
    \frac{4}{L^2}
    \frac{
        \left(e^{L\sqrt v/2}-1\right)^4
    }{
        e^{L\sqrt v/2}
    }.
\]
For standard self-concordance, \(L=2\), this becomes
\[
    F_2(v)
    =
    \frac{(e^{\sqrt v}-1)^4}{e^{\sqrt v}},
\]
recovering the self-concordant estimate above.




\end{document}
